\documentclass[11pt,a4paper]{article}
\usepackage[T1]{fontenc}
\usepackage[utf8]{inputenc}
\usepackage[spanish,es-tabla]{babel}
\usepackage{lmodern}
\usepackage{geometry}
\geometry{left=2.1cm,right=2.1cm,top=2.1cm,bottom=2.1cm}
\usepackage{array,longtable,booktabs,tabularx,xcolor}
\usepackage{microtype}
\usepackage{hyperref}
\hypersetup{colorlinks=true,urlcolor=blue,linkcolor=black}
\usepackage{enumitem}
\setlength{\parskip}{5pt}
\setlength{\parindent}{0pt}
\renewcommand{\arraystretch}{1.22}
\definecolor{uteqverde}{RGB}{0,108,62}

% Campos de la matriz: ID, tipo, producto/proceso, sistema/software, característica ISO, justificación.
\newcommand{\req}[6]{%
\texttt{#1} & #2 & #3 & #4 & #5 & #6 \\ \hline
}
\begin{document}
\begin{center}
{\small UNIVERSIDAD TÉCNICA ESTATAL DE QUEVEDO}\\
{\small Facultad de Ciencias de la Computación y Diseño Digital -- Carrera de Ingeniería de Software}\\[5pt]
{\Large\bfseries FORM-TL-01: Identificación y clasificación de tipos de requisitos}\\[3pt]
{\normalsize Ingeniería de Requerimientos (ISR-401) -- 4.\textordmasculine{} nivel, paralelo A}
\end{center}
\begin{tabularx}{\textwidth}{@{}>{\bfseries}lX@{}}
Caso: & Plataforma de gestión de citas médicas\\
Estudiante: & Vaca Romero David Octavio\\
Docente: & PhD. Gleiston Cicerón Guerrero Ulloa\\
Fecha de entrega (SGA): & 7 de noviembre de 2026\\
\end{tabularx}

\section{Objetivo y alcance}
Clasificar los 20 enunciados del caso de estudio en las dimensiones especificadas por la guía del taller: tipo de requisito, producto/proceso y sistema/software. Cuando corresponda, identificar la característica y subcaracterística de calidad según ISO/IEC 25010:2023~\cite{iso25010}. Fundamentar las decisiones mediante definiciones de ingeniería de requisitos~\cite{iso29148}.

\section{Matriz de clasificación (E01--E20)}
La siguiente matriz es una propuesta argumentada que debe contrastarse con el material de la asignatura. Los enunciados originales se incluyen como anexo para mantener la trazabilidad. En los requisitos de proceso, el nivel \emph{sistema/software} se interpreta como el \emph{ámbito de trabajo}; no describe una función del producto. E09 se clasifica como funcional porque ordena registrar un estado; su origen es una regla de dominio. En E02 se toma como frontera de la solución el software que emite el recordatorio, no el servicio de telecomunicaciones.

{\footnotesize
\begin{longtable}{|p{0.9cm}|p{1.9cm}|p{1.7cm}|p{1.6cm}|p{2.7cm}|p{5.2cm}|}
\hline
\textbf{ID} & \textbf{Tipo} & \textbf{Producto / Proceso} & \textbf{Sistema / Software} & \textbf{ISO 25010} & \textbf{Justificación breve}\\ \hline
\endfirsthead
\multicolumn{6}{l}{\small Continuación de la matriz de clasificación}\\
\hline
\textbf{ID} & \textbf{Tipo} & \textbf{Producto / Proceso} & \textbf{Sistema / Software} & \textbf{ISO 25010} & \textbf{Justificación breve}\\ \hline
\endhead
% FORM-TL-01 - Clasificacion individual. No entregar con celdas pendientes.
% Sintaxis: \req{ID}{tipo (funcional/calidad/restriccion)}{producto/proceso}{sistema/software}{ISO 25010 o No aplica}{justificacion}
\req{E01}{Pendiente}{Pendiente}{Pendiente}{Pendiente}{Pendiente de justificar.}
\req{E02}{Pendiente}{Pendiente}{Pendiente}{Pendiente}{Pendiente de justificar.}
\req{E03}{Pendiente}{Pendiente}{Pendiente}{Pendiente}{Pendiente de justificar.}
\req{E04}{Pendiente}{Pendiente}{Pendiente}{Pendiente}{Pendiente de justificar.}
\req{E05}{Pendiente}{Pendiente}{Pendiente}{Pendiente}{Pendiente de justificar.}
\req{E06}{Pendiente}{Pendiente}{Pendiente}{Pendiente}{Pendiente de justificar.}
\req{E07}{Pendiente}{Pendiente}{Pendiente}{Pendiente}{Pendiente de justificar.}
\req{E08}{Pendiente}{Pendiente}{Pendiente}{Pendiente}{Pendiente de justificar.}
\req{E09}{Pendiente}{Pendiente}{Pendiente}{Pendiente}{Pendiente de justificar.}
\req{E10}{Pendiente}{Pendiente}{Pendiente}{Pendiente}{Pendiente de justificar.}
\req{E11}{Pendiente}{Pendiente}{Pendiente}{Pendiente}{Pendiente de justificar.}
\req{E12}{Pendiente}{Pendiente}{Pendiente}{Pendiente}{Pendiente de justificar.}
\req{E13}{Pendiente}{Pendiente}{Pendiente}{Pendiente}{Pendiente de justificar.}
\req{E14}{Pendiente}{Pendiente}{Pendiente}{Pendiente}{Pendiente de justificar.}
\req{E15}{Pendiente}{Pendiente}{Pendiente}{Pendiente}{Pendiente de justificar.}
\req{E16}{Pendiente}{Pendiente}{Pendiente}{Pendiente}{Pendiente de justificar.}
\req{E17}{Pendiente}{Pendiente}{Pendiente}{Pendiente}{Pendiente de justificar.}
\req{E18}{Pendiente}{Pendiente}{Pendiente}{Pendiente}{Pendiente de justificar.}
\req{E19}{Pendiente}{Pendiente}{Pendiente}{Pendiente}{Pendiente de justificar.}
\req{E20}{Pendiente}{Pendiente}{Pendiente}{Pendiente}{Pendiente de justificar.}

\end{longtable}
}

\section{Cuantificación de los cinco requisitos vagos}
% Umbrales propuestos para validación con las partes interesadas; NO son resultados medidos.
Se reformulan los cinco enunciados vagos. Los valores son \textbf{propuestas de aceptación que deben negociarse}, no valores ya acordados ni mediciones reales. Se toma como referencia el pico estimado de 300 usuarios simultáneos y el uso frecuente de Android de gama media/baja descritos en el caso.

{\footnotesize
\begin{longtable}{|p{1.0cm}|p{2.5cm}|p{6.4cm}|p{3.6cm}|}
\hline
\textbf{ID} & \textbf{Calidad ISO/IEC 25010:2023} & \textbf{Redacción medible propuesta} & \textbf{Verificación prevista}\\ \hline
\endfirsthead
\hline \textbf{ID} & \textbf{Característica} & \textbf{Criterio cuantificado} & \textbf{Prueba prevista}\\ \hline
\endhead
E03 & Eficiencia de desempeño / comportamiento temporal & Con 300 sesiones concurrentes, al menos el 95\% de consultas de horarios disponibles responderá en 3 segundos o menos durante una prueba de 30 minutos. & Prueba de carga: tiempo por solicitud en segundos; calcular percentil 95. \\ \hline
E07 & Seguridad / confidencialidad & El sistema impedirá la lectura de datos de salud por usuarios sin un permiso autorizado en el 100\% de un conjunto de 100 pruebas de acceso denegado por rol. & Ejecutar 100 intentos de consulta sin autorización; contar accesos indebidos (meta: 0). \\ \hline
E10 & Capacidad de interacción / inclusividad & Al menos 17 de 20 participantes de 60 años o más completarán sin ayuda una reserva en 5 minutos o menos utilizando Android 10 o superior de gama media/baja. & Prueba con participantes reales: registrar éxito, ayuda requerida y duración en minutos. \\ \hline
E14 & Fiabilidad / disponibilidad & La plataforma estará operativa al menos el 99,5\% del tiempo programado de atención (lunes a sábado, 07:00--19:00), calculado en cada mes calendario. & Monitorización mensual: minutos disponibles / minutos programados $\times 100$. \\ \hline
E18 & Compatibilidad / interoperabilidad & En un entorno de pruebas conectado con la API documentada del proveedor de facturación, al menos 99 de 100 solicitudes válidas de registro de cobro producirán una confirmación correcta en 5 segundos o menos. & Prueba de 100 integraciones válidas; verificar respuesta, coincidencia de datos y tiempo en segundos. \\ \hline
\end{longtable}
}

\textbf{Límites de estas propuestas:} el caso no establece acuerdos de nivel de servicio, permisos detallados ni diseño de API. Por tanto, esos elementos necesitan validación con el director médico, el área de TI y el proveedor. La prueba con adultos mayores no debe simular participantes o resultados.


\section{Comparación documental de criterios y discrepancias}
La guía solicita comparar las clasificaciones con un compañero y registrar las discrepancias observadas. \textbf{No se deben inventar discusiones ni consensos.} Complete esta sección tras una revisión cruzada real.

{\small
\begin{longtable}{|p{1.3cm}|p{3.5cm}|p{3.5cm}|p{4.8cm}|}
\hline
\textbf{ID} & \textbf{Primera propuesta} & \textbf{Propuesta alternativa} & \textbf{Resolución argumentada}\\ \hline
Por registrar & Pendiente. & Pendiente. & Pendiente.\\ \hline
\end{longtable}
}


\section{Uso de herramientas de inteligencia artificial}
ChatGPT (GPT-6) se utilizó como apoyo para estructurar LaTeX, redactar propuestas de clasificación y métricas, comparar documentalmente los archivos aportados y verificar la compilación. Se contrastaron los enunciados del caso y las matrices disponibles; las decisiones argumentadas y los umbrales siguen sujetos a revisión académica y validación de interesados. No se atribuyen conversaciones, acuerdos, pruebas de campo o resultados inexistentes.

\begin{thebibliography}{9}
\bibitem{iso29148} ISO/IEC/IEEE, \emph{ISO/IEC/IEEE 29148:2018, Systems and software engineering --- Life cycle processes --- Requirements engineering}, 2018.
\bibitem{iso25010} ISO/IEC, \emph{ISO/IEC 25010:2023, Systems and software engineering --- Systems and software Quality Requirements and Evaluation (SQuaRE) --- Product quality model}, 2023.
\bibitem{angelo2026} A. M. Espinoza Pivaque, "FORM-TL-01: Taller de tipos de requisitos", repositorio académico \emph{ISR401-EspinozaAngelo}, GitHub, 2026. [En línea]. Disponible: \url{https://github.com/Espinoza-A-1006/ISR401-EspinozaAngelo}. Consultado el 7 de noviembre de 2026.
\end{thebibliography}

\appendix
\section{Enunciados originales del caso}
\textbf{E01.} El paciente podrá reservar una cita eligiendo especialidad, médico, fecha y hora entre los horarios disponibles.\par
\textbf{E02.} La plataforma enviará al paciente un recordatorio por mensaje de texto 24 horas antes de su cita.\par
\textbf{E03.} La aplicación debe ser rápida.\par
\textbf{E04.} El recepcionista podrá registrar la llegada del paciente que se acerca a la ventanilla.\par
\textbf{E05.} La plataforma se desarrollará en Java 21 con base de datos PostgreSQL, tecnologías que domina el área de TI.\par
\textbf{E06.} El médico podrá consultar la lista de pacientes del día y marcar cada consulta como atendida.\par
\textbf{E07.} Los datos de salud de los pacientes deben estar seguros.\par
\textbf{E08.} La aplicación móvil debe funcionar en teléfonos con Android 10 o superior.\par
\textbf{E09.} Una cita cancelada con menos de cuatro horas de anticipación se registrará como inasistencia.\par
\textbf{E10.} La plataforma debe ser fácil de usar para los pacientes adultos mayores.\par
\textbf{E11.} El proyecto debe entregarse en un plazo máximo de seis meses.\par
\textbf{E12.} La sesión de un usuario se cerrará automáticamente tras 10 minutos sin actividad.\par
\textbf{E13.} El quiosco de la sala de espera leerá el número de cédula ingresado en su pantalla táctil e imprimirá un ticket con el número de turno.\par
\textbf{E14.} La plataforma debe estar siempre disponible.\par
\textbf{E15.} El director médico podrá generar un reporte mensual de inasistencias por especialidad.\par
\textbf{E16.} Cada incremento se validará con el director médico al final de cada iteración de dos semanas.\par
\textbf{E17.} El tratamiento de los datos personales de los pacientes debe cumplir la Ley Orgánica de Protección de Datos Personales del Ecuador.\par
\textbf{E18.} La plataforma debe integrarse sin problemas con el sistema de facturación actual.\par
\textbf{E19.} Si el servidor principal falla, la plataforma debe restablecer el servicio en menos de 30 minutos sin perder ninguna cita confirmada.\par
\textbf{E20.} Todo cambio en el código fuente será revisado por otro integrante del equipo antes de integrarse a la rama principal.\par

\end{document}
